\documentclass[11pt]{article}
\usepackage[T1]{fontenc}
\usepackage{amsmath,amssymb,amsthm,mathtools}
\usepackage[margin=1in]{geometry}
\newtheorem{theorem}{Theorem}
\newtheorem{lemma}{Lemma}
\newtheorem{corollary}{Corollary}
\newcommand{\bits}{\{-1,1\}}
\newcommand{\HH}{\mathsf H}
\newcommand{\II}{\mathsf I}
\newcommand{\E}{\mathbb E}
\newcommand{\wh}[1]{\widehat{#1}}
\begin{document}

\section*{A fiber-sparse conditional phase theorem}

All physical variables below are independent uniform signs and all entropies are in bits.
Let $X,Y$ be disjoint coordinate sets and let
$f:\bits^X\times\bits^Y\to\bits$ be Boolean.  Write
\[
 a_{R,T}=\wh f(R,T),\qquad
 p_X(R)=\sum_{T\subseteq Y}a_{R,T}^2,
\]
and let $(\mathcal S_X,\mathcal S_Y)$ have law
$\Pr[(\mathcal S_X,\mathcal S_Y)=(R,T)]=a_{R,T}^2$.
For $R$ with $p_X(R)>0$, define the physical Fourier fiber
\[
 \mathcal F_R=\{T\subseteq Y:a_{R,T}\ne0\}.
\]

\begin{theorem}[conditional phase transfer with sparse fibers]\label{thm:main}
Assume that every nonzero $a_{R,T}$ has $|R|\le d$, where $d\ge1$ is an integer,
and that $|\mathcal F_R|\le K$ for every $R$ with $p_X(R)>0$.
For a physical slice $f_y(x)=f(x,y)$, write
\[
 c_R(y)=\E_x f(x,y)\chi_R(x)=2^{1-d}m_R(y),\qquad m_R(y)\in\mathbb Z,
\]
and define
\[
 \Gamma_d(f_y)=4^{1-d}\sum_R m_R(y)^2\log_2(m_R(y)^2),
\]
with zero summands omitted. Then
\begin{equation}\label{eq:main}
 \HH(f)\le {2\over\ln2}\II_Y(f)+2(d-1)+\log_2K
          -\E_y\Gamma_d(f_y).
\end{equation}
In particular, if $\operatorname{Var}(f)\ge\nu>0$, then
\begin{equation}\label{eq:fei}
 \HH(f)\le
 \left({2\over\ln2}+{2(d-1)+\log_2K\over\nu}\right)\II(f)
 -\E_y\Gamma_d(f_y).
\end{equation}
Thus fixed $(d,K,\nu)$ gives a dimension-independent FEI theorem for this genuine
Boolean class.  If $f$ is balanced, one may take $\nu=1$.
\end{theorem}

\begin{proof}
The exact spectral chain rule is
\begin{equation}\label{eq:chain}
 \HH(f)=\HH(p_X)+\HH(\mathcal S_Y\mid\mathcal S_X).
\end{equation}
The second term is retained: conditionally on $\mathcal S_X=R$ its law is supported
on $\mathcal F_R$, so
\begin{equation}\label{eq:cond}
 \HH(\mathcal S_Y\mid\mathcal S_X)
 \le\sum_Rp_X(R)\log_2|\mathcal F_R|\le\log_2K.
\end{equation}
No part of the marginal influence payment is reused in \eqref{eq:cond}.

For completeness, the refined marginal estimate is derived next.  The hybrid law
\[
 W(y,R)=2^{-|Y|}c_R(y)^2
\]
has $R$-marginal $p_X$, conditional law $c_R(y)^2$ given $y$, and hence
\[
 \HH(p_X)=\E_y\HH(c(y)^2)+I_W(R;y).
\]
Because every Boolean slice has Fourier degree at most $d$, Boolean Fourier
granularity gives $c_R(y)=2^{1-d}m_R(y)$ with integral $m_R(y)$ and
$\sum_Rm_R(y)^2=4^{d-1}$.  Direct substitution gives the exact identity
\[
 \HH(c(y)^2)=2(d-1)-\Gamma_d(f_y).
\]
Moreover
\[
 I_W(R;y)=\sum_R\operatorname{Ent}_y(c_R^2)
 \le {2\over\ln2}\sum_R\E_y\sum_{j\in Y}(D_jc_R)^2
 ={2\over\ln2}\II_Y(f),
\]
by the uniform-cube logarithmic Sobolev inequality and Parseval in $X$.
Combining this marginal bound with \eqref{eq:chain}--\eqref{eq:cond} proves
\eqref{eq:main}.

Finally, the Boolean Poincar\'e inequality gives
$\II(f)\ge\operatorname{Var}(f)\ge\nu$, while $\II_Y(f)\le\II(f)$.
Therefore $2(d-1)+\log_2K\le[2(d-1)+\log_2K]\II(f)/\nu$,
which proves \eqref{eq:fei}.
\end{proof}

\begin{lemma}[exact tensor-admission law]\label{lem:tensor}
Let $g:\bits^Z\to\bits$ be Boolean on coordinates disjoint from $X\cup Y$, and put
$F(x,y,z)=f(x,y)g(z)$.  Then the $X$-marginal of $F$ equals that of $f$, the
$X$-degree bound is unchanged, and for every $R$ with $p_X(R)>0$,
\[
 \mathcal F_R(F)=\mathcal F_R(f)\times\operatorname{supp}(\wh g),\qquad
 |\mathcal F_R(F)|=|\mathcal F_R(f)|\,|\operatorname{supp}(\wh g)|.
\]
Also
\[
 \HH_F(\mathcal S_{Y\cup Z}\mid\mathcal S_X)
 =\HH_f(\mathcal S_Y\mid\mathcal S_X)+\HH(g).
\]
\end{lemma}

\begin{proof}
On disjoint blocks, $\wh F(R,T,U)=\wh f(R,T)\wh g(U)$.  The support identity and
cardinality multiplication follow, because a product is nonzero exactly when both
factors are nonzero.  Summing $\wh g(U)^2=1$ preserves $p_X$.  The conditional law
given $R$ is the product of the old conditional law and $P_g(U)=\wh g(U)^2$, so
conditional entropy is additive.
\end{proof}

Lemma~\ref{lem:tensor} is the required arbitrary-tensor guard.  Starting from any
fixed cyclic-marginal seed and tensoring by $g^{\otimes t}$ preserves that marginal
but multiplies each nonempty fiber size by
$|\operatorname{supp}(\wh g)|^t$.  If $\HH(g)>0$, then this support size is at least
$2^t$, so no fixed-$K$ class admits the hardness-preserving tensor sequence.  If the
support size is one, $g$ is a Fourier character and $\HH(g)=0$, so it carries no
spectral-entropy hardness.  This exclusion uses only exact Fourier support
factorization, not FEI.

The condition is stable under adjoining dummy variables, permuting coordinates
within $X$ or $Y$, and multiplying by a $Y$-character (which merely translates each
fiber).  Under a disjoint product the parameters obey the exact rules
$d=d_1+d_2$ when both factors use $X$-coordinates and $K=K_1K_2$ for the corresponding
$Y$-fibers.  No unproved restriction-closure assertion is made.

\section*{Scope}
This is a proper-class theorem and not full FEI.  The universal problem is precisely
that no dimension-independent $d,K,\nu$ are known for arbitrary Boolean functions,
nor is there a universal decomposition into such pieces with all costs paid.  The
result allows arbitrary correlations among the coordinates of the conditional
spectral law on $Y$; it assumes no product structure there.  It does not allow a
correlated physical input law.

\end{document}
